\documentclass[10pt,a4paper]{amsart}
\usepackage[margin=19mm]{geometry}
\usepackage{amsmath,amssymb,mathtools}
\usepackage[scr=rsfs,cal=euler]{mathalfa}
\usepackage{xfrac}
\usepackage[unicode,hidelinks,bookmarks=false]{hyperref}
\newcommand{\ie}{i.e.\ }
\newcommand{\cf}{cf.\ }
\newcommand{\resp}{respectively\ }
\newcommand{\Subsection}{\subsection}
\providecommand{\nobreakdash}{\nobreak\hbox{-}}
\setlength{\emergencystretch}{2em}
\multlinegap0pt
\usepackage{amssymb}
\usepackage{xfrac}
\newtheorem{thm}{Theorem}
\newcommand{\BNORM}[1]{\Bigg\lVert#1\Bigg\rVert}
\newcommand{\MGd}[2]{\Big\{ \, {#1}\;\Big|\; {#2} \,\Big\} }
\newcommand{\Mgd}[2]{\big\{\,  {#1}\;\big|\; {#2} \,\big\} }
\newcommand {\N}{\mathbb N}
\newcommand{\NORM}[1]{\Big\lVert#1\Big\rVert}
\newcommand{\Norm}[1]{\big\lVert#1\big\rVert}
\newcommand {\Q}{\mathbb Q}
\newcommand {\R}{\mathbb R}
\newcommand {\del}{ \; \big| \;}
\newcommand{\equi}{\Leftrightarrow}
\newcommand{\maps}[1]{\mathop{\overset{#1}\longrightarrow}}
\newcommand{\maths}[1]{\[\begin{split}{#1}\end{split}\]}
\newcommand{\norm}[1]{\lVert#1\rVert}
\newcommand{\om}{\omega}
\newcommand{\tom} {\emptyset}
\title{A Classification of Order Convergence via a Transfinite Fatou Hierarchy}
\author{Antonio Avil\'es, Christian Rosendal, Mitchell A. Taylor, Pedro Tradacete}
\date{Source: arXiv:2604.02588v1 (CC BY 4.0)}
\begin{document}
\maketitle

\noindent\emph{Source and licence.} A Classification of Order Convergence via a Transfinite Fatou Hierarchy. Version arXiv:2604.02588v1 (CC BY 4.0), distributed by arXiv under the Creative Commons Attribution 4.0 International licence, http://creativecommons.org/licenses/by/4.0/, as recorded in the arXiv licence metadata for this exact version and shown on the abstract page https://arxiv.org/abs/2604.02588v1. Full licence notice: \texttt{licenses/arxiv-2604.02588-CC-BY-4.0.txt}.\par\smallskip

\noindent\textit{Editorial scope.} Counted unit: the article's thm thm (source0.tex lines 713--715). The article's complete original proof of it is delivered with it (source0.tex lines 717--948, one \texttt{proof} environment of 232 lines). No mathematics is written, repaired or re-derived: the statement and the proof are the article's own lines, in the article's own notation, and no retained line is substituted or re-flowed.  No further source-local context is needed: every cross-reference of every retained line resolves inside the two delivered spans.  Nothing was omitted from any retained span, so the package carries no omission record.  No retained line contains a citation, so the wrapper carries no bibliography and no bibliography entry.  No retained line contains a figure environment, an image-inclusion command or drawing code, so no image is delivered and the package declares no asset dependency.  Editorial formatting only, in addition: an AMS \texttt{amsart} preamble that restates the article's own declarations and macros used by the retained lines, namely the environments \texttt{thm} on separate counters, together with the standard packages those lines need.  The retained environments print in this document as Theorem 1 in order of appearance, whereas the article prints the same items with the numbering of its own full document.  The complete native source of the article as distributed in the arXiv e-print for this exact version is bundled unchanged as \texttt{sources/197748/source0.tex}.
\begin{thm} 
For every countable ordinal $\alpha$, there is a separable Banach lattice with a countable $\pi$-basis that fails to be $\alpha$-Fatou.
\end{thm}

\begin{proof}
Our proof goes by induction on $1\leqslant \alpha<\om_1$. For each $\alpha$, we will construct several objects,
\begin{enumerate}
    \item a separable Banach lattice $(X_\alpha,\norm{\cdot}_\alpha)$,
    \item a countable $\pi$-basis $B_\alpha$ for $X_\alpha$,
    \item a Banach lattice homomorphism $X_\alpha\maps{\phi_\alpha} \mathbb{R}$ of norm $1$,
    \item a sequence 
$$
z^\alpha_1\geqslant z^\alpha_2\geqslant  \ldots> 0=\inf_nz^\alpha_n
$$
whose terms belong to the set 
$$
S_\alpha=\big\{x\in X^+_\alpha\del \norm{x}_\alpha=\phi_\alpha(x)=1\big\},
$$
\item\label{last property} a countable tree $T_\alpha\subseteq \Psi\big((z^\alpha_n)\big)\cap S_\alpha^{<\N}$ so that
$$
\rho(T_\alpha)>\alpha.
$$
\end{enumerate}
Observe then that, by property \eqref{last property}, we have
$$
\rho\Big(\Psi\big((z^\alpha_n)\big)\Big)>{\alpha}
$$
and hence that $X_\alpha$ fails to be $\alpha$-Fatou.



\


{\bf Base case, $\alpha=1$.} We let $X_1$ be the Banach lattice  $c$ of convergent real sequences endowed with the equivalent renorming 
$$
\Norm{(t_1,t_2,\ldots)}_1 =\max\Big\{ \tfrac{1}{3}\Norm{(t_1,t_2,\ldots)}_{\ell_\infty},  \big|\lim_nt_n\big|\Big\}.
$$ 
Note that 
$$
B_1=\Mgd{(0,\ldots, 0,t,0,\ldots)}{t\in \Q_+}
$$
is a countable $\pi$-basis for $X_\alpha$.
Define $\phi_1$ by
$$
\phi_1\big((t_1,t_2,\ldots)\big) = \lim_n t_n
$$
and, for each $n$, set 
$$
z^1_n = (\underbrace{0,0,\ldots,0}_{n\text{ times}}, 1,1,1,\ldots) \in S_1.
$$ 
Finally, letting $y_1 = (1,1,1,\ldots)\in S_1$, we see that $(y_1)\in \Psi\big((z^\alpha_n)\big)$ and therefore that
$$
T_1=\big\{\tom, (y_1)\big\}\subseteq \Psi\big((z^\alpha_n)\big)\cap S_1^{<\N}
$$
and that $\rho(T_1)=2>1$.

\

{\bf Successor case.} Suppose that $(X_\alpha,\norm{\cdot}_\alpha)$, $B_\alpha$, $\phi_\alpha$, $(z_n^\alpha)$ and $T_\alpha$ have been defined as above. We then let $X_{\alpha+1}=X_\alpha$ with the new norm 
$$
\norm x_{\alpha+1} = \max\Big\{\tfrac17{\norm x_\alpha},\big|\phi_\alpha(x)\big|\Big\}.
$$ 
Observe that $\norm\cdot_{\alpha+1}\leqslant \norm\cdot_\alpha$, whereby $\norm{\phi_\alpha}_{\alpha+1}=1$ and we may therefore set 
 $\phi_{\alpha+1} =\phi_\alpha$ and $B_{\alpha+1}=B_\alpha$.
 Note also that
$$
\norm x_{\alpha+1}=\norm x_\alpha \;\equi\;  \norm x_\alpha=|\phi_\alpha(x)|,
$$
from which it follows that $S_{\alpha}\subseteq S_{\alpha+1}$. We may thus set $z^{\alpha+1}_n = z^\alpha_n\in S_\alpha\subseteq S_{\alpha+1}$ for all $n$.
Finally,  let 
$$
T_{\alpha+1}
=\big\{ (z_1^\alpha,y_1,\ldots, y_n)\del (y_1,\ldots, y_n)\in T_\alpha\big\}\cup\{\tom\} \;\;\subseteq \;\; S_{\alpha}^{<\N} \;\;\subseteq \;\; S_{\alpha+1}^{<\N}
$$
and note that $\rho(T_{\alpha+1})>\rho(T_{\alpha})>\alpha$.
To see that $T_{\alpha+1}\subseteq  \Psi\big((z^{\alpha+1}_n)\big)$, note first that, for all $x,y\in S_{\alpha}$, we have 
$$
\phi_\alpha\big((x-y)^+\big)
=\big(\phi_\alpha(x-y)\big)^+
=0^+
=0
=\phi_\alpha(x-y),
$$
whereby $\Norm{(x-y)^+}_{\alpha+1}=\tfrac17\Norm{(x-y)^+}_{\alpha}$ and $\norm{x-y}_{\alpha+1}=\tfrac17\norm{x-y}_{\alpha}$.
Thus, for all $n$, 
\maths{
\Norm{(z_1^{\alpha} - z^{\alpha+1}_n)^+}_{\alpha+1} 
&=\Norm{(z_1^{\alpha} - z^{\alpha}_n)^+}_{\alpha+1} \\
&=\tfrac{1}{7}\Norm{(z_1^{\alpha} - z^{\alpha}_n)^+}_\alpha\\
&\leqslant\tfrac{1}{7}\norm{z^\alpha_1}_\alpha\\
& =\tfrac{1}{7}\norm{z^\alpha_1}_{\alpha+1}\\
& < \frac{\norm{z_1^{\alpha}}_{\alpha+1}}{3^1}
}
and 
\maths{
\Norm{(y_i - z^{\alpha+1}_n)^+}_{\alpha+1} 
&=\Norm{(y_i - z^{\alpha}_n)^+}_{\alpha+1} \\
%&= \max\left\{\frac{\Norm{(y_i - z^{\alpha+1}_n)^+}_\alpha}{4},\Big|\phi_\alpha\big({(y_i - z^{\alpha+1}_n)^+}\big)\Big|\right\}\\
%&= \tfrac{1}{5}\Norm{(y_{i}-z^{\alpha+1}_n)^+}_\alpha\\
&= \tfrac{1}{7}\Norm{(y_{i}-z^\alpha_n)^+}_\alpha\\
&\leqslant  \frac{1}{7}\frac{\|{y}_{i}\|_\alpha}{3^{i}} \\
&<  \frac{\norm{y_i}_{\alpha+1}}{3^{i+1}}.
} 
Similarly, 
\maths{
\norm{y_1 -z_1^{\alpha} }_{\alpha+1}
& = \tfrac17\Norm{y_1-z^\alpha_1}_\alpha \leqslant \frac{2}{7}<  \frac{\norm{z_1^\alpha}_{\alpha+1}}{3^1}
}
and 
\maths{
\norm{{y}_{i+1} - {y}_i}_{\alpha+1} 
= \tfrac17\Norm{y_{i+1}-y_{i}}_\alpha
\leqslant \frac17\frac{\norm{y_i}_\alpha}{3^{i}} 
< \frac{\norm{y_i}_{\alpha+1}}{3^{i+1}}.
} 
Together, these inequalities establish that $(z_1^\alpha,y_1,\ldots, y_n)\in \Psi\big((z^{\alpha+1}_n)\big)$ whenever $(y_1,\ldots, y_n)\in T_\alpha$ and hence that $T_{\alpha+1}\subseteq \Psi\big((z^{\alpha+1}_n)\big)$.

	
	\
	
{\bf Limit case.} Suppose that $\alpha_1 <\alpha_2 < \ldots<\alpha = \lim_n \alpha_n$. Assume also that the construction has been done for all ordinals smaller than $\alpha$. We first consider the $\ell_\infty$-sum of the Banach lattices $X_{\alpha_n}$,
$$
X^\infty_\alpha 
= \bigoplus_{\ell_\infty}\MGd{X_{\alpha_n} }{ n\in \N} 
= \MGd{x=(x^1,x^2,\ldots)\in \prod_{n\in \N}X_{\alpha_n} }
{ \sup_n\|x^n\|_{\alpha_n}<\infty}.
$$
Endowed with coordinatewise operations and the norm $\|x\|_\alpha = \sup_n \|x^n\|_{\alpha_n}$, this is a Banach lattice. Furthermore, because each $\phi_{\alpha_n}$ is a Banach lattice homomorphism of norm 1, it follows that 
$$
X^\phi_\alpha =\MGd{x=(x^1,x^2,\ldots)\in X^\infty_\alpha }{ \lim_n \phi_{\alpha_n}(x^n) \text{ exists}}
$$ 
is a Banach sublattice of $X^\infty_\alpha$ and that $\phi\colon X^\phi_\alpha\to \R$ defined by
$$
\phi\big((x^n)\big)=\lim_n\phi_{\alpha_n}(x^n)
$$
is a lattice homomorphism of norm $1$. Observe also that $$X^\infty_\alpha\cap \prod_n S_{\alpha_n}\subseteq X^\phi_\alpha.
$$
Unfortunately, there is no reason for $X^\phi_\alpha$ to be separable and so $X_\alpha$ will be chosen to be a specific  separable Banach sublattice of $X^\phi_\alpha$. 


For this we first define the uncountable tree
$T^\phi_\alpha$ of all finite strings of elements of $X^\phi_\alpha$ of the form 
$$
\Big(
\big(\underbrace{0,0,\ldots,0}_{k-1},{y_1^k},{y_1^{k+1}},{y_1^{k+2}},\ldots\big)
\,,\;\; \ldots\;\; , \, 
\big(\underbrace{0,0,\ldots,0}_{k-1},{y_n^k},{y_n^{k+1}},{y_n^{k+2}},\ldots\big)
\Big),
$$
where $k\geqslant 1$ and $(y_1^m,y_2^m,\ldots, y_n^m)\in T_{\alpha_m}$ for all $m\geqslant k$. We observe that $\rho(T^\phi_\alpha)\geqslant \alpha+1.$ 
This is because the rank of an element of the tree as above is given by 
$$
\min\Mgd{\rho_{T_{\alpha_m}}\big((y_1^m,y_2^m,\ldots, y_n^m)\big) } { m\geqslant k},
$$ 
and therefore, by the inductive hypothesis, we can find strings of rank at least $\alpha_m$ for every $m$, and the rank of the root $\emptyset$ is $\alpha$. 

By induction on countable ordinals $\beta$, it is easily seen that, if a tree $\Upsilon$ satisfies $\rho(\Upsilon)\geqslant \beta$, then there is a countable subtree $\Upsilon'\subseteq \Upsilon$ with $\rho(\Upsilon')\geqslant \beta$. 
Applying this to $\Upsilon=T^\phi_\alpha$, we take a countable tree $T_\alpha\subseteq T^\phi_\alpha$ such that $\rho(T_\alpha)>\alpha$.

We define $X_\alpha$ to be the separable Banach sublattice of $X^\phi_\alpha$ generated by 
\begin{enumerate}
    \item all elements of $X^\phi_\alpha$ that appear in the tree $T_\alpha$, 
    \item all eventually zero vectors $(x_1,\ldots,x_n,0,0,\ldots)\in X^\phi_\alpha$,
    \item vectors $z^\alpha_k = \left({z^{\alpha_1}_k},{z^{\alpha_2}_k},{z^{\alpha_3}_k},\ldots\right)$ for  $k<\omega$.  
\end{enumerate}


If $B_{\alpha_n}$ is the countable $\pi$-basis for $X_{\alpha_n}$, then the collection $B_\alpha$ of vectors of the form $(0,\ldots,0,b,0,0,\ldots)$
with $b\in B_{\alpha_n}$ in the $n$th position form a countable $\pi$-basis for $X_\alpha$. The sequence $z^\alpha_k = \left({z^{\alpha_1}_k},{z^{\alpha_2}_k},{z^{\alpha_3}_k},\ldots\right)$ that we already defined is contained in $S_\alpha$ and satisfies 
$$
z^\alpha_1\geqslant z^\alpha_2\geqslant \ldots> 0=\inf_kz^\alpha_k.
$$

The nodes of the tree $T_\alpha$ indeed belong to $S_\alpha^{<\mathbb{N}}$ and by construction we have that $\rho(T_\alpha)>\alpha$. It thus remains to show that $T_{\alpha}\subseteq  \Psi\big((z^{\alpha}_m)\big)$. 
First,
\maths{
\BNORM{
\Big(\underbrace{0,0,\ldots,0}_{k-1}&,{y_{i+1}^k},y_{i+1}^{k+1},{y_{i+1}^{k+2}},\ldots\Big)
-
\Big(\underbrace{0,0,\ldots,0}_{k-1},{y_i^k},{y_i^{k+1}},{y_i^{k+2}},\ldots\Big)
}_\alpha
\\
&
=\BNORM{
\Big(\underbrace{0,0,\ldots,0}_{k-1},  \big(  {y_{i+1}^k-y_i^k}\big),   \big(  {y_{i+1}^{k+1}-y_i^{k+1}}\big),    \big(y_{i+1}^{k+2}-y_i^{k+2}\big),\ldots\Big)
}_\alpha
\\
&
=\sup_{m\geqslant k}\NORM{ {y_{i+1}^m-y_i^m} }_{\alpha_m}
\\
&
\leqslant \sup_{m\geqslant k}\tfrac 1{3^i}\Norm{ {y_i^m} }_{\alpha_m}
\\
&
\leqslant \frac 1{3^i}
\NORM{
\Big(\underbrace{0,0,\ldots,0}_{k-1},{y_{i}^k},{y_{i}^{k+1}},{y_{i}^{k+2}},\ldots\Big)
}_\alpha.
}
Similarly, for all $k$ and $m$,
\maths{
\BNORM{
\Bigg(&\Big(\underbrace{0,0,\ldots,0}_{k-1},{y_{i}^k},{y_{i}^{k+1}},{y_{i}^{k+2}},\ldots\Big)
-
z^\alpha_m\Bigg)^+
}_\alpha
\\
&
=\BNORM{
\Bigg(\Big(\underbrace{0,0,\ldots,0}_{k-1},{y_{i}^k},{y_{i}^{k+1}},{y_{i}^{k+2}},\ldots\Big)
-
\Big({z^{\alpha_1}_m},{z^{\alpha_2}_m},{z^{\alpha_3}_m},\ldots\Big)\Bigg)^+
}_\alpha
\\
&
=\BNORM{\Big(
\underbrace{0,0,\ldots,0}_{k-1},
{\big({y_{i}^k}-{z^{\alpha_k}_m}\big)^+},
{\big({y_{i}^{k+1}}-{z^{\alpha_{k+1}}_{m}}\big)^+},{\big({y_{i}^{k+2}}-{z^{\alpha_{k+2}}_{m}}\big)^+},\ldots\Big)
}_\alpha
\\
&
=\sup_{r\geqslant k}    \NORM{        \big(    y_{i}^r      -     z^{\alpha_r}_m     \big)^+            }_{\alpha_r} 
\\
&
\leqslant \sup_{r\geqslant k} \tfrac1{3^i}\Norm{y_{i}^r}_{\alpha_r} 
\\
&
\leqslant \frac 1{3^i}
\NORM{
\Big(\underbrace{0,0,\ldots,0}_{k-1},{y_{i}^k},{y_{i}^{k+1}},{y_{i}^{k+2}},\ldots\Big)
}_\alpha.
} 
This completes the verification of properties (1)-(5) and hence the inductive step.
\end{proof}

\end{document}
